% !TEX TS-program = xelatex
% !TEX encoding = UTF-8 Unicode

\documentclass{resume}
\usepackage{zh_CN-Adobefonts_external}
\usepackage{linespacing_fix}
\usepackage{cite}

\begin{document}
\pagenumbering{gobble}

\name{张钧瑞}

\contactInfo{(+86) xxx-xxxx-xxxx}{your.email@example.com}{应用统计 硕士研究生}{GitHub @Zjrua}
 
\section{个人总结}
哈尔滨工业大学应用统计硕士研究生，本科信息与计算科学专业。具有扎实的数学建模与统计分析基础，熟练掌握 Python 与 PyTorch 进行数据处理与建模。多次参加数学建模竞赛并获奖，目前在导师指导下从事体质数据统计分析研究，论文已投稿《体育科学》。高中学过信息学竞赛，具备良好的编程与算法基础。

\section{教育背景}
\datedsubsection{\textbf{哈尔滨工业大学}, 应用统计, \textit{硕士研究生}}{2025.9 -- 2028.6}
\ 研究方向为统计建模与数据分析，参与体质数据相关课题研究
\datedsubsection{\textbf{哈尔滨工业大学}, 信息与计算科学, \textit{理学学士}}{2020.9 -- 2024.6}
\ 数学与计算机交叉学科背景

\section{技术能力}
\begin{itemize}[parsep=0.2ex]
  \item \textbf{编程语言}: Python, C/C++, SQL, LaTeX
  \item \textbf{框架与库}: PyTorch, pandas, NumPy, scikit-learn
  \item \textbf{工具与平台}: Linux/macOS, Git, Jupyter, VS Code
\end{itemize}

\section{科研经历}
\datedsubsection{\textbf{体质数据分析研究 | 导师课题组}}{2025.09 -- 至今}
\begin{itemize}
  \item 负责体质健康数据的清洗、统计建模与分析全流程
  \item 应用多元统计方法探索影响体质的关键因素
  \item 撰写学术论文并投稿《体育科学》
\end{itemize}

\section{竞赛获奖}
\datedsubsection{\textbf{全国统计建模大赛}（研究生组）}{2026}
\ 国家级竞赛，参赛编号 TJJM20260414180979
\datedsubsection{\textbf{美国大学生数学建模竞赛 (MCM/ICM)}}{2022}
\ \textbf{Meritorious Winner}（一等奖，全球前 7\%）
\datedsubsection{\textbf{全国大学生数学建模竞赛 (CUMCM)}}{2023}
\ \textbf{黑龙江省二等奖}

\section{其他}
\begin{itemize}[parsep=0.2ex]
  \item GitHub: \href{https://github.com/Zjrua}{https://github.com/Zjrua}
  \item 高中阶段参加全国信息学奥林匹克竞赛 (NOI)
  \item 入党积极分子
\end{itemize}

\end{document}
